% !TEX root = ../../main.tex

\section{约束分析}

在系统需求分析过程中，除功能性需求和非功能性需求外，还需要明确系统开发和运行过程中受到的约束条件。
约束分析有助于限定系统实现边界，避免在毕业设计阶段将系统目标扩大为完整的开放域手语翻译系统，也有助于后续系统设计和测试时形成更合理的评价标准。
结合 HearBridge 系统的实际开发条件和应用场景，本文主要从开发周期、技术平台、数据资源、模型能力、运行环境和应用边界等方面进行约束分析。

\subsection{开发周期约束}

本系统作为本科毕业设计项目，开发周期相对有限，需要在规定时间内完成需求分析、系统设计、功能实现、联调测试和论文撰写等工作。
HearBridge 系统同时涉及 HarmonyOS 移动端、Spring Boot 服务端、Vue 后台管理端和 Python 手势识别服务，系统组成较多，各端之间还需要进行接口联调和功能验证。
因此，系统需求必须控制在可实现、可测试和可展示的范围内，避免引入过多超出周期承载能力的复杂功能。

在该约束下，系统将重点放在手语资源浏览、实时手势训练、句子视频识别、样本管理、模型训练和模型版本发布等核心功能上，而没有将目标设定为大规模开放域连续手语翻译系统。
对于连续手语识别、复杂语义翻译、多模态对话和完整数字人自动生成等功能，本文将其作为后续扩展方向，而不是当前系统必须完成的功能范围。

\subsection{技术平台约束}

本系统移动端基于 HarmonyOS 平台开发，客户端功能实现需要遵循 HarmonyOS 应用模型、权限机制、页面路由方式和多媒体访问规则。
相机采集、视频选择、网络访问和页面状态管理等能力均需要在 HarmonyOS 平台提供的能力范围内完成。
因此，移动端功能设计需要兼顾平台特性和实际设备运行条件，不能完全按照普通 Web 应用或 Android 原生应用的方式设计交互流程。

服务端采用 Spring Boot 技术栈，后台管理端采用 Vue 3 与 Element Plus，识别服务采用 Python 生态中的 FastAPI、MediaPipe 和 TensorFlow/Keras 等工具。
各端技术栈不同，数据格式、接口调用方式和异常处理方式也存在差异。
系统需要通过统一接口约定和清晰的数据结构降低跨端协作成本，同时避免某一端承担过多不适合自身技术环境的职责。
例如，复杂视觉识别和模型训练不适合直接放在移动端完成，因此系统将其拆分到 Python 识别服务中实现。

\subsection{数据资源约束}

手势识别和句子视频识别效果依赖样本数据的规模、质量和覆盖范围。
由于本系统开发周期有限，难以构建大规模、多用户、多场景和多角度的完整手语数据集。
当前系统更适合在固定词表、固定训练项和受限句子场景下进行识别能力验证，而不适合直接支持开放词表或复杂连续手语表达。

此外，不同国家和地区的手语体系存在差异，公开数据集中的手语类别、动作习惯和标注方式不一定能直接对应系统中的训练资源。
对于自采样本而言，样本质量又会受到采集设备、光照条件、拍摄角度、动作规范程度和采集人数等因素影响。
因此，系统需要提供样本管理、样本质量查看和模型版本管理等能力，以便后续逐步扩展数据规模并优化模型效果。

\subsection{模型能力约束}

HearBridge 系统中的识别模型主要用于固定词表下的手势动作识别和有限词表句子视频识别。
该模型能够为用户训练和辅助交流提供一定识别参考，但不能保证在所有真实场景下都获得稳定准确的识别结果。
当用户动作不规范、视频画面模糊、手部遮挡严重、光照条件较差或动作类别相似度较高时，识别结果可能出现漏识别、误识别或置信度较低等问题。

句子视频识别功能同样受到模型能力约束。
系统可以对受限词表下的短视频进行词级识别和语义修正，但这并不等同于完整的连续手语识别或开放域手语翻译。
语义修正功能只能在原始识别候选结果基础上进行整理和自然化表达，不能替代视觉识别模型本身。
如果原始识别结果错误较多，语义修正结果也可能受到影响。
因此，系统在展示识别结果时需要保留原始识别信息和修正结果，避免将语义修正结果误认为绝对正确的翻译结论。

\subsection{运行环境约束}

系统采用多端协同架构，移动端、服务端、后台管理端、对象存储服务和 Python 识别服务之间需要通过网络进行通信。
实际运行效果会受到网络状态、服务部署地址、接口可用性、文件访问权限和服务器性能等因素影响。
若识别服务不可用、对象存储访问失败或网络延迟较高，系统中的实时识别、视频上传识别和资源展示功能都会受到影响。

实时手势训练对相机帧传输和识别反馈的连续性有一定要求，因此设备性能、摄像头质量、网络连接和识别服务处理速度都会影响用户体验。
句子视频识别则涉及视频文件选择、上传、分析和结果返回，处理耗时通常高于普通业务接口，因此系统需要提供等待提示和异常提示，避免用户在处理过程中重复提交或误判系统无响应。

\subsection{应用边界约束}

HearBridge 系统定位为面向听障辅助训练和受限场景辅助交流的工程原型，主要用于验证移动端训练应用、后台资源管理和手势识别服务协同工作的可行性。
系统可以为用户提供手语动作示范、实时训练反馈和句子视频识别结果，但其识别结果仍属于辅助性信息，不能替代专业手语教师、医疗康复评估或正式无障碍翻译服务。

在实际使用中，系统应避免对识别能力进行过度承诺。
对于训练反馈，系统主要帮助用户了解动作是否接近目标动作；
对于句子视频识别，系统主要验证有限词表条件下从视频到词序列再到自然语言展示的基本链路。
后续若要扩展到真实公共服务场景，还需要进一步扩大样本规模、提高模型泛化能力、完善用户评估机制，并结合听障用户反馈持续优化系统交互和识别结果展示方式。
